\scanpage{19}
\mspage{--- 47 ---}
\begin{center}
\textbf{Algèbre et géométrie}

M. Jean-Pierre \textsc{Serre}, professeur
\end{center}

Soit $\mathrm{K}$ un corps complet pour une valuation discrète, de corps résiduel $k$. On suppose que $\mathrm{K}$ est de caractéristique zéro et que $k$ est algébriquement clos de caractéristique $p$. Soit $\bar{\mathrm{K}}$ une clôture algébrique de $\mathrm{K}$, et soit $\mathrm{G}$ le groupe de Galois de $\bar{\mathrm{K}}$ sur $\mathrm{K}$. Soit $\mathrm{V}$ un espace vectoriel de dimension finie sur le corps $p$-adique $\mathbf{Q}_p$, et soit
\[
\varrho:\mathrm{G}\longrightarrow\operatorname{Aut}(\mathrm{V})
\]
un homomorphisme continu. On dit alors que $\mathrm{V}$ (ou $\varrho$) est une \textit{représentation $p$-adique} de $\mathrm{G}$.

Soit $\chi$ l’homomorphisme de $\mathrm{G}$ dans le groupe des unités $p$-adiques correspondant à l’action naturelle de $\mathrm{G}$ sur les racines de l’unité d’ordre une puissance de $p$. Soit $\mathrm{C}$ le complété de $\bar{\mathrm{K}}$. Si $\mathrm{V}$ est, comme ci-dessus, une représentation $p$-adique de $\mathrm{G}$, posons $\mathrm{V}_{\mathrm{C}}=\mathrm{C}\otimes\mathrm{V}$ (le produit tensoriel étant pris sur $\mathbf{Q}_p$) ; pour tout $n\in\mathbf{Z}$, soit $\mathrm{V}_{\mathrm{C}}^n$ le $\mathrm{K}$-sous-espace de $\mathrm{V}_{\mathrm{C}}$ formé des éléments $x$ tels que $\varrho(s)x=\chi(s)^nx$ pour tout $s\in\mathrm{G}$, et soit $\mathrm{V}_{\mathrm{C}}(n)=\mathrm{C}\otimes_{\mathrm{K}}\mathrm{V}_{\mathrm{C}}^n$. D’après un résultat de \textsc{Tate}, la somme directe des $\mathrm{V}_{\mathrm{C}}(n)$ s’identifie à un sous-espace vectoriel de $\mathrm{V}_{\mathrm{C}}$. On dit que $\mathrm{V}$ admet une \textit{décomposition de Hodge $p$-adique} (ou que $\mathrm{V}$ est un \textit{module de Hodge}) si ce sous-espace vectoriel est égal à $\mathrm{V}_{\mathrm{C}}$, autrement dit si l’on a :
\[
\mathrm{V}_{\mathrm{C}}=\sum_{n\in\mathbf{Z}}\mathrm{V}_{\mathrm{C}}(n).
\]

Soit $\mathrm{V}$ un tel module, soit $\mathrm{G}_{\mathrm{V}}=\varrho(\mathrm{G})$ le sous-groupe correspondant de $\operatorname{Aut}(\mathrm{V})$, et soit $\mathrm{H}_{\mathrm{V}}$ \textit{l’enveloppe algébrique} de $\mathrm{G}_{\mathrm{V}}$, autrement dit le plus petit sous-groupe algébrique (sur $\mathbf{Q}_p$) du groupe linéaire $\mathbf{GL}_{\mathrm{V}}$ dont le groupe des points contienne $\mathrm{G}_{\mathrm{V}}$.

Le sujet du cours a été l’étude du groupe $\mathrm{H}_{\mathrm{V}}$, et en particulier la démonstration des résultats suivants :
% archive sheet number: 18
